\documentclass[conference]{IEEEtran}
\usepackage{amsmath,amssymb}
\usepackage{booktabs}
\usepackage{array}

\begin{document}

\title{A Minimal IEEEtran Dual-Column Hello World}

\author{\IEEEauthorblockN{Toolchain Verification}
\IEEEauthorblockA{\textit{medvlthinker-imgdiff-compute} \\
Build test for \texttt{build\_ieee.sh}}}

\maketitle

\begin{abstract}
This one-page document confirms that the IEEEtran conference class compiles to a
real two-column PDF with the local toolchain. It exercises the abstract, a table,
and inline mathematics such as the confidence margin threshold $\tau$ and a symbol
$\phi$, plus a display equation.
\end{abstract}

\begin{IEEEkeywords}
IEEEtran, toolchain, tectonic, two-column
\end{IEEEkeywords}

\section{Introduction}
This is column one of a two-column IEEE conference layout. If you can read this
text flowing into a second column on the same page, the dual-column format is
working. The gate escalates to the strong model whenever the margin falls below
the threshold $\tau$, i.e. $m < \tau$. A companion symbol $\phi$ is used below in
a display equation to confirm math-mode rendering across both inline and display
contexts. Lorem ipsum filler is added so the text is long enough to wrap from the
first column into the second one and demonstrate the balanced two-column body.

\subsection{A Display Equation}
\begin{equation}
\phi(\tau) = \sum_{i=1}^{N} \mathbb{1}\!\left[m_i < \tau\right] \cdot c_i,
\end{equation}
where $c_i$ is the per-sample cost and $m_i$ the confidence margin.

\section{A Table}
Table~\ref{tab:demo} shows a minimal booktabs table rendered inside the
IEEE column.

\begin{table}[t]
\centering
\caption{Cascade operating point (illustrative).}
\label{tab:demo}
\begin{tabular}{@{}lcc@{}}
\toprule
Model & Accuracy & Rel. FLOPs \\
\midrule
7B (cap320)      & 0.512 & 0.18 \\
Cascade ($\tau{=}0.426$) & 0.572 & 0.74 \\
32B (think)      & 0.572 & 1.00 \\
\bottomrule
\end{tabular}
\end{table}

\section{Conclusion}
If this PDF has two columns, a title block, an abstract, a table, and the symbols
$\tau$ and $\phi$ render correctly, the IEEEtran build path is verified end to end.

\end{document}
